\documentclass[10pt,a4paper]{amsart}
\usepackage[margin=19mm]{geometry}
\usepackage{amsmath,amssymb,mathtools}
\usepackage[scr=rsfs,cal=euler]{mathalfa}
\usepackage{xfrac}
\usepackage[unicode,hidelinks,bookmarks=false]{hyperref}
\newcommand{\ie}{i.e.\ }
\newcommand{\cf}{cf.\ }
\newcommand{\resp}{respectively\ }
\newcommand{\Subsection}{\subsection}
\providecommand{\nobreakdash}{\nobreak\hbox{-}}
\setlength{\emergencystretch}{2em}
\multlinegap0pt
\usepackage{amssymb}
\usepackage{breqn}
\usepackage{xcolor}
\usepackage{mathtools}
\usepackage{dsfont}
\newtheorem{theorem}{Theorem}
\newtheorem{prop}{Proposition}
\newtheorem{lemma}{Lemma}
\usepackage{cleveref}
\crefname{theorem}{Theorem}{Theorems}
\crefname{prop}{Proposition}{Propositions}
\crefname{lemma}{Lemma}{Lemmas}
\newcommand{\C}{\mathbb{C}}
\newcommand{\Q}{\mathbb{Q}}
\newcommand{\R}{\mathbb{R}}
\newcommand{\Z}{\mathbb{Z}}
\newcommand{\wt}{\widetilde}
\title{Lifting all elements in $\mathrm{SL}_n(\mathbb{Z}/q\mathbb{Z})$}
\author{Amitay Kamber, P\'eter P. Varj\'u}
\date{Source: arXiv:2310.10269v2 (CC BY 4.0)}
\begin{document}
\maketitle

\noindent\emph{Source and licence.} Lifting all elements in $\mathrm{SL}_n(\mathbb{Z}/q\mathbb{Z})$. Version arXiv:2310.10269v2 (CC BY 4.0), distributed by arXiv under the Creative Commons Attribution 4.0 International licence, http://creativecommons.org/licenses/by/4.0/, as recorded in the arXiv licence metadata for this exact version and shown on the abstract page https://arxiv.org/abs/2310.10269v2. Full licence notice: \texttt{licenses/arxiv-2310.10269-CC-BY-4.0.txt}.\par\smallskip

\noindent\textit{Editorial scope.} Counted unit: the article's theorem roots (source0.tex lines 277--281). The article's complete original proof of it is delivered with it (source0.tex lines 630--687, one \texttt{proof} environment of 58 lines, which the article places away from the statement). No mathematics is written, repaired or re-derived: the statement and the proof are the article's own lines, in the article's own notation, and no retained line is substituted or re-flowed.  Retained uncounted as the source-local context the proof needs: lines 501--504; lines 530--537; lines 573--590; lines 606--625.  Nothing was omitted from any retained span, so the package carries no omission record.  No retained line contains a citation, so the wrapper carries no bibliography and no bibliography entry.  No retained line contains a figure environment, an image-inclusion command or drawing code, so no image is delivered and the package declares no asset dependency.  Editorial formatting only, in addition: an AMS \texttt{amsart} preamble that restates the article's own declarations and macros used by the retained lines, namely the environments \texttt{theorem}, \texttt{prop}, \texttt{lemma} on separate counters, together with the standard packages those lines need.  The retained environments print in this document as Theorem 1, Proposition 1, Lemma 1 in order of appearance, whereas the article prints the same items with the numbering of its own full document.  The complete native source of the article as distributed in the arXiv e-print for this exact version is bundled unchanged as \texttt{sources/148531/source0.tex}.
\begin{prop}\label{prop:small P factor}
Assume that $q$ has a small $P$ factor. Then there is $\beta\in(\Z/q\Z)^{\times}$
such that $\beta^{n}=1$ and $\left|n\beta\right|>q^{1-1/k}-n$. 
\end{prop}

\begin{lemma}
\label{lem:Enough residues}For every $m$, there is $C>0$ such that the following holds.
Let $q\in \Z_{>C}$ be without small P factors.
Then there
are $\alpha_{1},\ldots,\alpha_{m}\in(\Z/q\Z)^{\times}$ such that each
$\alpha_{i}$ has an $n$-th root $\beta_{i}\in(\Z/q\Z)^\times$, (i.e., $\beta_{i}^{n}=\alpha_{i}$),
$0<\wt\alpha_{i}<C$ and  $(\wt\alpha_{i}/\wt\alpha_{j})^{1/n}\notin\Q$ for $i\ne j$.
\end{lemma}

\begin{lemma}[Bohr set lemma]
\label{lem:Bohr Set Lemma}
For each $k\in\Z_{>0}$ and $c>0$, there is a constant $C$ such that 
the following holds.
Let $A\subset\Z/q\Z$ be such that $\gcd(A\cup\{q\})=1$
and let $\rho<C^{-1}q^{-1/(k+1)}$.
Then there is $C^{-1}\rho<\rho'\le\rho$,
$r\le k$ and $u_{1},\ldots,u_{r}\in\Z/q\Z$, $l_{1},\ldots,l_{r}\in\Z_{>0}$
such that
\begin{enumerate}
\item $|u_{i}a|\le c\rho'q$ for all $a\in A$, i.e., $u_{i}\in B(A,c\rho')$,
\item $B(A,\rho')\subset\left\{ \sum_{i=1}^{r}n_{i}u_{i}:|n_{i}|\le l_{i}\right\} \subset B(A,C\rho')$,
\item if $\left|n_{i}\right|\le C^{-1}l_{i}/\rho'$ and $\sum_{i=1}^{r}n_{i}u_{i}\in B(A,t\rho')$
for some $1\le t\le{\rho'}^{-1}/2$, then $\left|n_{i}\right|\le Ctl_{i}$, and
\item the map $(n_{1},\ldots,n_{r})\to\sum_{i=1}^{r}n_{i}u_{i}$ is injective
for $\left|n_{i}\right|\le C^{-1}l_{i}/\rho'$.
\end{enumerate}
\end{lemma}

\begin{lemma}\label{lem:lifting beta}
For all $D_1,D_2\in \R_{>0}$, there is $C\in \R_{>0}$ such that the following holds.
Let $A\subset\Z/q\Z$ such that $\gcd(A\cup\{q\})=1$, and let $\alpha,\beta\in\Z/q\Z$
such that $\beta A\subset A\cup \alpha A$ and $|\alpha|\le D_1$.
Let $\rho\in(0,C^{-1})$, $r\in\Z_{>0}$, $u_1,\ldots,u_r\in\Z/q\Z$
and $l_1,\ldots, l_r\in\Z_{>0}$ be such that
\begin{enumerate}
\item $u_{i}\in B(A,D_1^{-1}\rho)$,
\item $B(A,\rho)\subset\left\{ \sum_{i=1}^{r}n_{i}u_{i}:|n_{i}|\le l_{i}\right\} \subset B(A,D_2\rho)$, and
\item If $\left|n_{i}\right|\le D_2^{-1}l_{i}/\rho$ and $\sum_{i=1}^{r}n_{i}u_{i}\in B(A,t\rho)$
for some $1\le t\le{\rho}^{-1}/2$, then $\left|n_{i}\right|\le D_2tl_{i}$.
\end{enumerate}

Then there is a matrix $B=(B_{i,j})\in M_r(\Z)$ such that
\[
\beta (u_1,\ldots, u_r)\equiv B(u_1,\ldots,u_r) \mod q
\]
and $|B_{i,j}|\le C l_j/l_i$ for all $i$, $j$.
(Recall that the notation $(u_1,\ldots,u_r)$ is understood as a column vector.)
\end{lemma} 

\begin{theorem}\label{thm:large roots}
	For every $k>0,n>0$, there is a constant $C$ such that the following holds.
	For every $q>0$, there are $\alpha,\beta\in(\Z/q\Z)^\times$ with
	$|\alpha|<C$, $\beta^{n}=\alpha$ and $|n\beta|>C^{-1}q^{1-1/k}$.
\end{theorem}

\begin{proof}[Proof of \cref{thm:large roots}]
By \cref{prop:small P factor}, we may assume that $q$ has no small $P$-factors. 
	
Let $m=n^{k}+1$ (the reason for this choice will be explained right at the end of the proof), and let $C_1$,
$\alpha_{1},\ldots,\alpha_{m}$ be as in \cref{lem:Enough residues}, i.e., $0<\wt\alpha_{s}<C_1$. Let $\beta_{s}$ be some $n$-th root of $\alpha_{s}$. 

Choose $c=C_1^{-1}$ in \cref{lem:Bohr Set Lemma}, and let $C_2$ be the constant
$C$ in that lemma.
	
Consider the set 
\[
A=\left\{ \beta_{1}^{a_{1}}\cdot\ldots\cdot\beta_{m}^{a_{m}}:0\le a_{s}<n\right\} .
\]
Let $\rho=q^{1/k}$. 
We suppose to the contrary that the theorem is false.
As we already discussed, this implies that $\lvert n a\rvert \le C_2^{-1}\rho q$ for all $a\in A$,
and hence $n\in B(A,\rho')$ for all $\rho'>C^{-1}\rho$.
In particular, we get $r\ge 1$ when we apply \cref{lem:Bohr Set Lemma}
for $A$ and $\rho$.


Note that $\rho<C_2^{-2}q^{-1/(k+1)}$ if $q$ is sufficiently large.
By \cref{lem:Bohr Set Lemma}, there is $C_2^{-1}\rho<\rho'\le\rho$, $0<r\le k$,
$u_{1},\ldots,u_{r}\in \Z/q\Z$ and $l_{1},\ldots,l_{r}\in \Z_{>0}$ such that 
items (1)--(4) in the lemma hold.
We can now apply \cref{lem:lifting beta} for each pair $\alpha_s, \beta_s$ in the
role of $\alpha$ and $\beta$ and with $\rho'$ in the role of $\rho$, $C_1$ in the role of $D_1$
and $\max(C_1,C_2)$
in the role of $D_2$.
Denoting by $C_3$ the constant $C$ in \cref{lem:lifting beta}, we get
matrices $B^{(s)} \in M_n(\Z)$ with
\[
\beta_s(u_1,\ldots, u_r)\equiv B^{(s)}(u_1,\ldots, u_r) \mod q
\]
and $|B^{(s)}_{ij}|\le C_3l_j/l_i$.

We observe that
\begin{align*}
B^{(s_1)}B^{(s_2)}(u_1,\ldots, u_r)&\equiv\beta_{s_1}\beta_{s_2}(u_1,\ldots, u_r) \\
&\equiv B^{(s_2)}B^{(s_1)}(u_1,\ldots, u_r) \mod q.    
\end{align*}
The $i,j$ entries of both matrices $B^{(s_1)}B^{(s_2)}$ and $B^{(s_2)}B^{(s_1)}$
are bounded by $rC_3^2l_j/l_i\le C_2^{-1} l_j/\rho'$ provided $q$ is sufficiently large.
By the injectivity property in item (4) of \cref{lem:Bohr Set Lemma},
we get $B^{(s_1)}B^{(s_2)}=B^{(s_2)}B^{(s_1)}$.
A very similar argument gives $(B^{(s)})^n=\wt\alpha_sI$.

	
We finally get $m=n^{k}+1$ commuting matrices $B^{(1)},\ldots,B^{(n^{k}+1)}\in M_{r}(\Z)$, $r\le k$, such that $(B^{(s)})^{n}=\wt\alpha_{s}I$. 
Since the minimal polynomial of $B^{(s)}$ has no double roots, 
$B^{(s)}$ is diagonalizable, and since all the matrices commute, 
there is a matrix $P \in M_r(\C)$ such that all the matrices $PB^{(s)}P^{-1}$ are diagonal. 
Moreover, the diagonal of $PB^{(s)}P^{-1}$ is of the form $\wt\alpha_{s}^{1/n}(\zeta_{n}^{t_{s1}},\ldots,\zeta_{n}^{t_{sr}})$,
for $\wt\alpha_{s}^{1/n}\in\R_{>0}$, $\zeta_{n}=\exp(2\pi i/n)$ and $0\le t_{sj}\le n-1$. 
By the pigeon-hole principle, there are $s_{1}\neq s_{2}$ with $t_{s_{1}j}=t_{s_{2}j}$ for all $j$. 
This implies that $B^{(s_{1})}$ is a scalar multiple of $B^{(s_{2})}$, that is, $B^{(s_{1})}=\left(\wt\alpha_{s_1}/\wt\alpha_{s_2}\right)^{1/n}B^{(s_2)}$.
This is a contradiction since those are integer matrices and $\left(\wt\alpha_{s_1}/\wt\alpha_{s_2}\right)^{1/n}\notin\Q$.
\end{proof}

\end{document}
